% =====================================================================
%  Bölüm 17 — Projeler (Cilt I)
% =====================================================================
\chapter{Projeler}\label{bol:proje}

\begin{bolumhedefleri}
\begin{itemize}
  \item Bir formalizasyon projesini ifadeden temiz bir ispata kadar
        yürütmenin yöntemini,
  \item Cilt~I'in konularını derinleştiren on iki projeyi ve her biri için
        önerilen bir yol haritasını,
  \item Projeler arasındaki bağımlılıkları ve Cilt~II'ye uzanan köprüleri.
\end{itemize}
\end{bolumhedefleri}

Bir alıştırma birkaç satırlık bir ispatla biter. Bir \emph{proje} ise
günler, bazen haftalar sürer. Önce doğru tanımları bulmak, sonra ispatı
önsavlara bölmek ve en sonunda tanımlarınızın ispatı ne kadar
kolaylaştırdığını ya da zorlaştırdığını görmek gerekir. Kitap boyunca bazı
sonuçları ``bir projedir'' diyerek kâğıt üzerinde bıraktık. Church--Rosser
teoremi, güçlü normalleşme ve sözdizimsel kesme eleme bunlardandır. Bu
bölüm, bunları on iki proje olarak bir araya getiriyor. Her projenin hedef
teoremi eşlik eden \Lean{} projesinde \lc{sorry} ile hazır bekliyor:
\lc{ModernMantik/}\allowbreak\lc{Alistirmalar/}\allowbreak\lc{Projeler.lean}.

% ---------------------------------------------------------------------
\section{Bir formalizasyon projesi nasıl yürütülür?}\label{sec:proje-yontem}

Deneyimli formalizasyoncuların neredeyse hepsi aynı döngüyü izler
(\cref{fig:proje-dongusu}). Bölüm~12'nin tamlık ispatını bu döngüyle
yazmıştık: Önce \lc{tamlik} ifadesi, sonra onu besleyen önsavlar, en son
da her önsavın kendi ispatı geldi.

\begin{figure}[htbp]
\centering
\begin{tikzpicture}[
    adim/.style={draw=#1, fill=#1!8, very thick, rounded corners=5pt,
                 text width=2.5cm, align=center, inner sep=5pt,
                 font=\small, minimum height=1.5cm},
    ok/.style={-{Stealth[length=2.6mm]}, very thick, black!40},
    no/.style={font=\baslikfont\scriptsize\bfseries, text=white,
               fill=#1, circle, inner sep=1.6pt},
  ]
  \node[adim=anarenk]        (a) at (0,0)    {\textbf{İfade}\\{\scriptsize tanımlar ve hedef teorem, \lc{sorry} ile}};
  \node[adim=ornekframe]     (b) at (3.4,0)  {\textbf{İskelet}\\{\scriptsize önsavlar, her biri \lc{sorry}}};
  \node[adim=ikincil]        (c) at (6.8,0)  {\textbf{İspat}\\{\scriptsize önsavları birer birer kapat}};
  \node[adim=teoremframe]    (d) at (10.2,0) {\textbf{Temizlik}\\{\scriptsize \leankw{\#print} \lc{axioms}, belgeler}};
  \draw[ok] (a) -- (b); \draw[ok] (b) -- (c); \draw[ok] (c) -- (d);
  \draw[ok, dikkatframe, dashed] (c.south) to[bend left=35]
       node[below, font=\baslikfont\scriptsize, text=dikkatframe]{ispat zorlaşıyor mu? tanımları değiştir} (a.south);
  \foreach \n/\c/\k in {a/anarenk/1, b/ornekframe/2, c/ikincil/3, d/teoremframe/4}
    \node[no=\c] at (\n.north west) {\k};
\end{tikzpicture}
\caption{Bir formalizasyon projesinin döngüsü. Kesikli ok en önemli
adımdır: Zor bir ispat çoğu zaman yanlış seçilmiş bir tanımın
işaretidir.}
\label{fig:proje-dongusu}
\end{figure}

\begin{onemlifikir}[Önemli Fikir: \lc{sorry} ile geliştirme]
\lc{sorry}, \Lean{}'e ``bu kısmı sonra ispatlayacağım'' demenin yoludur.
Bir projenin iskeletini önce \lc{sorry}'lerle kurup sonra her \lc{sorry}'yi
ayrı ayrı doldurmak, ispatın genel yapısının doğru olup olmadığını
erkenden görmenizi sağlar. İskeletteki bir önsav yanlış ifade edilmişse
(örneğin bir varsayım eksikse), bunu ana teoremi ispatlamaya çalışırken
fark edersiniz. Önsavı ispatlamak için saatler harcamadan önce.
\end{onemlifikir}

\noindent Kitaptaki tasarım tartışmalarından çıkan birkaç pratik kural:

\begin{itemize}
  \item \textbf{Hesaplayan tanımlar seçin.} \lc{deger}'i \lc{Bool} değerli
        tanımlamak (Bölüm~10), Kalmár önsavını \leantac{cases} ile satır satır
        ilerletmeyi mümkün kıldı.
  \item \textbf{Önermeyi güçlendirin.} Kalmár önsavı hem doğruyu hem yanlışı
        taşıdı. Zayıflatma ispatı \lc{Δ}'yı genelleştirdi. Tümevarım
        takıldığında ilk soru budur.
  \item \textbf{İndis aritmetiğine dikkat.} Bağımlı tiplerde \lc{n + 0} ile
        \lc{0 + n} farklıdır (Bölüm~15). Tanımı özyinelemenin yönüne göre
        yazın.
  \item \textbf{Aksiyomları denetleyin.} Yapıcı olması gereken bir sonuç
        için \leankw{\#print} \lc{axioms} son adımdır (Bölüm~16).
\end{itemize}

% ---------------------------------------------------------------------
\section{Projelerin haritası}\label{sec:proje-haritasi}

\Cref{fig:proje-haritasi}, on iki projeyi dayandıkları bölümlere ve
zorluklarına göre yerleştiriyor. Oklar projeler arasındaki
bağımlılıkları gösteriyor: Örneğin P6'yı yapmak için önce P4'ün ilk
yarısı ve P5'in fikirleri gerekir.

\begin{figure}[htbp]
\centering
\begin{tikzpicture}[
    proje/.style={draw=#1, fill=#1!12, thick, rounded corners=3pt,
                  font=\baslikfont\small\bfseries, inner sep=4pt, minimum width=1.0cm},
    ok/.style={-{Stealth[length=2mm]}, thick, black!40},
    x=1.75cm, y=1.25cm,
  ]
  % eksenler
  \foreach \b/\x in {10/0, 11/1, 12/2, 13/3, 14/4, 15/5, 16/6}
    \node[font=\baslikfont\scriptsize, text=black!55] at (\x,-0.75) {Bölüm \b};
  \draw[black!20] (-0.5,-0.45) -- (6.5,-0.45);
  \node[font=\baslikfont\scriptsize, text=ikincil, anchor=east]     at (-0.55,0) {\zorluk{2}};
  \node[font=\baslikfont\scriptsize, text=dikkatframe, anchor=east] at (-0.55,1.5) {\zorluk{3}};
  \node[font=\baslikfont\scriptsize, text=vurgu, anchor=east]       at (-0.55,3) {\zorluk{3}+};
  \foreach \y in {0,1.5,3} \draw[black!10, dashed] (-0.5,\y) -- (6.5,\y);
  % projeler
  \node[proje=ikincil]     (p1)  at (0,0)   {P1};
  \node[proje=dikkatframe] (p4)  at (1,1.5) {P4};
  \node[proje=dikkatframe] (p12) at (2,1.5) {P12};
  \node[proje=dikkatframe] (p2)  at (3,1.5) {P2};
  \node[proje=vurgu]       (p5)  at (3,3)   {P5};
  \node[proje=vurgu]       (p6)  at (2.3,3) {P6};
  \node[proje=ikincil]     (p3)  at (4,0)   {P3};
  \node[proje=dikkatframe] (p7)  at (4,1.5) {P7};
  \node[proje=vurgu]       (p8)  at (4.6,3) {P8};
  \node[proje=vurgu]       (p9)  at (3.8,3) {P9};
  \node[proje=ikincil]     (p10) at (5,0)   {P10};
  \node[proje=ikincil]     (p11) at (6,0)   {P11};
  % bağımlılıklar
  \draw[ok] (p1) -- (p12);
  \draw[ok] (p4) -- (p6);
  \draw[ok] (p5) -- (p6);
  \draw[ok] (p7) -- (p9);
  \draw[ok] (p3) -- (p8);
  \draw[ok, dashed, teoremframe] (p4.north) -- ++(0,0.65)
       node[above, font=\baslikfont\scriptsize, text=teoremframe]{Cilt~II: Kripke};
\end{tikzpicture}
\caption{On iki projenin haritası. Yatay eksen projenin dayandığı ana
bölümü, dikey eksen zorluğu gösterir. Oklar ``önce bunu yapın''
anlamındadır.}
\label{fig:proje-haritasi}
\end{figure}

\begin{table}[htbp]
\centering
\small
\renewcommand{\arraystretch}{1.35}
\rowcolors{2}{tablosatir}{white}
\begin{tabularx}{\linewidth}{>{\bfseries}p{0.8cm} X p{1.7cm} p{1.4cm} p{2.1cm}}
\rowcolor{tabloust}
\tablobaslik{} & \tablobaslik{Proje} & \tablobaslik{Bölüm} & \tablobaslik{Zorluk} & \tablobaslik{Emek} \\
P1  & Doğruluk tablosu yönteminin doğruluğu        & 10     & \zorluk{2} & 1--2 gün \\
P2  & Ardışık hesapta ispat aramanın doğruluğu     & 13     & \zorluk{3} & 3--5 gün \\
P3  & Normalleştiricinin doğruluğu                 & 14     & \zorluk{2} & 1--2 gün \\
P4  & Yapıcı doğal çıkarım ve Peirce yasası         & 4, 11  & \zorluk{3} & 1 hafta \\
P5  & Sözdizimsel kesme eleme                      & 13     & \zorluk{3} & 2--4 hafta \\
P6  & LJ ile yapıcı doğal çıkarımın denkliği       & 5, 13  & \zorluk{3} & 1--2 hafta \\
P7  & Tip korunumu                                 & 14     & \zorluk{3} & 1 hafta \\
P8  & Church--Rosser teoremi                       & 14     & \zorluk{3} & 1--2 hafta \\
P9  & Güçlü normalleşme                            & 14     & \zorluk{3} & 2--3 hafta \\
P10 & Biriktiricili ters çevirme                   & 15     & \zorluk{2} & 1 gün \\
P11 & Tamsayıların halka yapısı                    & 16     & \zorluk{2} & 2--3 gün \\
P12 & Yapıcı model varlığı                         & 10, 12 & \zorluk{3} & 1 hafta \\
\end{tabularx}
\caption{Projelerin özeti. Emek tahminleri, kitabı buraya kadar
alıştırmalarıyla birlikte izlemiş bir okur içindir.}
\label{tab:projeler}
\end{table}

% ---------------------------------------------------------------------
\section{A. Karar yöntemleri ve yansıma}\label{sec:projeler-a}

Bu gruptaki üç proje, kitapta yazdığımız \emph{çalışan} fonksiyonların
(\lc{totolojiMi}, \lc{ara}, \lc{adim?}) gerçekten söylediklerini yaptığını
ispatlamayı ister. Bir fonksiyonun doğruluğunu ispatlamak, onu bir
\emph{karar yöntemine} dönüştürmenin ilk adımıdır. Cilt~II'nin 6. bölümünde
bu tekniği, \emph{yansıma} adıyla, \Lean{}'in kendi önermelerine
uygulayacağız.

\begin{projekarti}{P1}{Doğruluk tablosu yöntemi doğrudur}
\projebilgi{Bölüm~10}{\zorluk{2}}{1--2 gün}
Bölüm~10'daki \lc{totolojiMi}, formülü \lc{tumDegerlemeler}'deki her
değerlemede hesaplar. Bunun gerçekten totolojiliği yakaladığını ispatlayın.

\leandosya[P1]{Alistirmalar/Projeler}{p1}

\textbf{Önerilen yol.}
\begin{enumerate}[nosep]
  \item \emph{Kapsama önsavı}: Her \lc{w} için, \lc{tumDegerlemeler n}
        içinde ilk \lc{n} atomda \lc{w} ile anlaşan bir \lc{v} vardır.
        \lc{n} üzerinde tümevarımla ve \lc{genislet}'in tanımıyla.
  \item Kapsama önsavını \lc{yerellik} ile birleştirin: \lc{deger v A =
        deger w A}.
  \item \lc{List.all} için \lc{List.all\_eq\_true} önsavını kullanın.
\end{enumerate}
\end{projekarti}

\begin{projekarti}{P2}{İspat arama doğrudur}
\projebilgi{Bölüm~13}{\zorluk{3}}{3--5 gün}
Bölüm~13'teki \lc{ara} fonksiyonu \lc{true} döndürdüğünde, ardışığın gerçekten
kesmesiz bir türetimi vardır.

\leandosya[P2]{Alistirmalar/Projeler}{p2}

\textbf{Önerilen yol.}
\begin{enumerate}[nosep]
  \item \lc{tam\_aux}'un örüntülerini aynen izleyin; her durumda \lc{ara}'nın
        tanımı tek bir özyinelemeli çağrıya ya da bir \lc{\&\&}'e açılır.
  \item \lc{(a \&\& b) = true}'yu \lc{Bool.and\_eq\_true} ile ikiye bölün.
  \item Atomik durum için \lc{List.any\_eq\_true} ve \lc{List.contains}
        önsavlarından ortak bir atom elde edin.
  \item \emph{Uzantı:} Ters yönü (\lc{Gecerli} ise \lc{ara = true}) de
        ispatlayıp \lc{ara}'yı bir karar yöntemine dönüştürün.
\end{enumerate}
\end{projekarti}

\begin{projekarti}{P3}{Normalleştiricinin her adımı bir β-adımıdır}
\projebilgi{Bölüm~14}{\zorluk{2}}{1--2 gün}
\leandosya[P3]{Alistirmalar/Projeler}{p3}

\textbf{Önerilen yol.}
\begin{enumerate}[nosep]
  \item \lc{t} üzerinde tümevarım yapın. \lc{app t₁ t₂} durumunda
        \lc{t₁}'in bir \lc{lam} olup olmadığına göre ayırın; çünkü
        \lc{adim?}'ın örüntüleri örtüşüyor.
  \item \lc{Option.map f x = some u} hipotezini \lc{x = some u'} ve \lc{u =
        f u'} biçimine çeviren küçük bir önsav yazın.
  \item \emph{Uzantı:} \lc{adim? t = none} ise \lc{Normal t}.
\end{enumerate}
\end{projekarti}

% ---------------------------------------------------------------------
\section{B. İspat kuramı}\label{sec:projeler-b}

\begin{projekarti}{P4}{Yapıcı doğal çıkarım ve Peirce yasası}
\projebilgi{Bölüm~4, 11}{\zorluk{3}}{1 hafta}
\lc{Turetilir}'den RAA kuralını çıkarınca yapıcı doğal çıkarım elde edilir
(Bölüm~11, \cref{sec:gecerlilik-sinir}). Peirce yasası bu sistemde
türetilemez. Ama bunu göstermek için iki değerli anlambilim yetmez.

\leandosya[P4]{Alistirmalar/Projeler}{p4}

\textbf{Önerilen yol.}
\begin{enumerate}[nosep]
  \item (a) kural tümevarımıyla birkaç satırdır.
  \item (b) için bir \emph{Kripke modeli} tanımlayın: kısmi sıralı bir
        dünyalar kümesi ve tekdüze bir atom yorumu. Zorlama bağıntısını
        ($w \Vdash A$) formül üzerinde özyinelemeyle tanımlayın (Bölüm~4).
  \item \lc{TuretilirY}'nin Kripke anlambilimine göre geçerliliğini ispatlayın.
        \lc{iseG} durumunda tekdüzelik önsavı gerekecek.
  \item Peirce yasası için iki dünyalı bir karşı model kurun:
        $w_0 \le w_1$, $p$ yalnızca $w_1$'de doğru.
\end{enumerate}
\emph{Kaynak:} \citet{kripke1965,troelstra1988}. Kripke anlambilimi
Cilt~II'nin 2. bölümünde ayrıntılı olarak işlenecek.
\end{projekarti}

\begin{projekarti}{P5}{Sözdizimsel kesme eleme}
\projebilgi{Bölüm~13}{\zorluk{3}}{2--4 hafta}
Bölüm~13'te kesmenin elenebilirliğini anlambilim yoluyla, birkaç satırda
ispatladık. Gentzen'in yolu ise kesmeli bir türetimi doğrudan kesmesiz bir
türetime \emph{dönüştürmektir}. Bu proje, bu dönüşümü yazmayı ister.

\leandosya[P5]{Alistirmalar/Projeler}{p5}

\textbf{Önerilen yol.}
\begin{enumerate}[nosep]
  \item Kesmesiz \lc{Ardisik} için \emph{sözdizimsel} tersinirlik önsavları
        ispatlayın (Bölüm~13'tekiler anlambilimseldi). Örneğin \lc{Ardisik
        ((A {\sembolkodfont ⋀} B) :: Γ) Δ → Ardisik (A :: B :: Γ) Δ}, türetim üzerinde
        tümevarımla.
  \item \emph{Kesme kabul edilebilirliği}: \lc{Ardisik Γ (A :: Δ)} ve
        \lc{Ardisik (A :: Γ) Δ} ise \lc{Ardisik Γ Δ}. \lc{A}'nın boyutu
        üzerinde dış tümevarım, birinci türetim üzerinde iç tümevarım.
        Ana formül her iki tarafta da yeni tanıtılmışsa tersinirlik
        önsavlarını kullanın.
  \item Son olarak \lc{ArdisikK} üzerinde tümevarımla, her kesmeyi bir
        önceki adımla eleyin.
\end{enumerate}
\emph{Kaynak:} \citet{gentzen1935,pfenning2000,negri2001}.
\end{projekarti}

\begin{projekarti}{P6}{LJ ile yapıcı doğal çıkarım aynıdır}
\projebilgi{Bölüm~5, 13}{\zorluk{3}}{1--2 hafta}
Gentzen'in sezgici ardışık hesabı LJ, sağ tarafta en fazla bir formüle izin
verir. LJ'nin, P4'teki \lc{TuretilirY} ile aynı formülleri türettiğini
gösterin.

\leandosya[P6]{Alistirmalar/Projeler}{p6}

\textbf{Önerilen yol.}
\begin{enumerate}[nosep]
  \item LJ $\Rightarrow$ doğal çıkarım: Sol kurallar eleme kurallarıyla
        taklit edilir. Kolay yöndür.
  \item Doğal çıkarım $\Rightarrow$ LJ: Eleme kuralları için LJ'de
        \emph{kesme} gerekir. Önce LJ için kesme kabul edilebilirliğini
        (P5'in tekniğiyle) ispatlayın.
  \item \lc{iseL} kuralının ana formülü öncülde de tuttuğuna dikkat edin.
        Bu, LJ'de daraltmanın kabul edilebilirliği için gereklidir.
\end{enumerate}
\emph{Kaynak:} \citet{negri2001,prawitz1965,dyckhoff1992}.
\end{projekarti}

% ---------------------------------------------------------------------
\section{C. $λ$-hesabı}\label{sec:projeler-c}

\begin{projekarti}{P7}{Tip korunumu}
\projebilgi{Bölüm~14}{\zorluk{3}}{1 hafta}
Bir terim indirgendiğinde tipi değişmez. Bu teorem, ``iyi tiplenmiş
programlar yanlış gitmez'' sloganının yarısıdır.

\leandosya[P7]{Alistirmalar/Projeler}{p7}

\textbf{Önerilen yol.}
\begin{enumerate}[nosep]
  \item \emph{Zayıflatma}: \lc{Tiplenir (Γ₁ ++ Γ₂) t A} ise
        \lc{Tiplenir (Γ₁ ++ B :: Γ₂) (kaydir 1 Γ₁.length t) A}.
        \lc{t} üzerinde tümevarım, \lc{Γ₁}'i genelleştirerek.
  \item \emph{Yerine koyma}: \lc{Tiplenir (Γ₁ ++ B :: Γ₂) t A} ve
        \lc{Tiplenir Γ₂ s B} ise \lc{Tiplenir (Γ₁ ++ Γ₂) (yerine
        Γ₁.length s t) A}. \lc{bul} için liste birleştirme önsavları
        gerekecek.
  \item Son olarak \lc{Adim} üzerinde tümevarım. Yalnızca \lc{beta}
        durumu yerine koyma önsavını kullanır.
\end{enumerate}
\emph{Kaynak:} \citet{barendregt1984,sorensen2006}.
\end{projekarti}

\begin{projekarti}{P8}{Church--Rosser teoremi}
\projebilgi{Bölüm~14}{\zorluk{3}}{1--2 hafta}
\leandosya[P8]{Alistirmalar/Projeler}{p8}

\textbf{Önerilen yol.}
\begin{enumerate}[nosep]
  \item \emph{Paralel indirgeme} \lc{Paralel t u}'yu tanımlayın: Bir
        terimdeki redekslerin herhangi bir alt kümesi aynı anda
        indirgenir.
  \item \lc{Adim ⊆ Paralel ⊆ Adimlar} olduğunu gösterin.
  \item Takahashi'nin yöntemi \citep{takahashi1995}: Her terimin bütün
        redekslerini indirgeyen \lc{tamGelisme t}'yi tanımlayın ve
        \lc{Paralel t u → Paralel u (tamGelisme t)} önsavını ispatlayın.
        Bu, \lc{Paralel}'in elmas özelliğini hemen verir.
  \item Elmas özelliğini \lc{Adimlar}'a taşıyın. \lc{kaydir} ve \lc{yerine}
        için bir ``yerine koyma ile paralel indirgeme'' önsavı gerekecek.
\end{enumerate}
\end{projekarti}

\begin{projekarti}{P9}{Güçlü normalleşme}
\projebilgi{Bölüm~14}{\zorluk{3}}{2--3 hafta}
Tiplenebilen her terimin her indirgeme dizisi sonludur. \lc{GucluNormal},
bu özelliğin \emph{tümevarımsal} bir tanımıdır: Bir terim, bütün tek adımlık
indirgenmişleri güçlü normalleşiyorsa güçlü normalleşir.

\leandosya[P9]{Alistirmalar/Projeler}{p9}

\textbf{Önerilen yol.}
\begin{enumerate}[nosep]
  \item Isınma: \lc{omega\_adim} ile \lc{p9\_omega}.
  \item Tait'in yöntemi: Tip üzerinde özyinelemeyle bir
        \lc{Indirgenebilir : Tip → Terim → Prop} tanımlayın. Ok tipinde
        ``her indirgenebilir argümana uygulandığında indirgenebilir''.
  \item Üç özellik (CR1--CR3): İndirgenebilir terimler güçlü normalleşir,
        indirgenme altında kapalıdır ve ``nötr'' terimler için genişleme
        altında kapalıdır.
  \item \emph{Temel önsav}: Tiplenebilen bir terime indirgenebilir terimler
        yerine konursa sonuç indirgenebilirdir. P7'nin yerine koyma
        önsavları burada yeniden kullanılır.
\end{enumerate}
\emph{Kaynak:} \citet{tait1967}; \citet[Bölüm~6]{girard1989}.
\end{projekarti}

% ---------------------------------------------------------------------
\section{D. Bağımlı tipler ve bölümler}\label{sec:projeler-d}

\begin{projekarti}{P10}{Biriktiricili ters çevirme}
\projebilgi{Bölüm~15}{\zorluk{2}}{1 gün}
\leandosya[P10]{Alistirmalar/Projeler}{p10}

\textbf{Önerilen yol.}
\begin{enumerate}[nosep]
  \item \lc{p10\_tersAux bos w}'nin tipi \lc{Vek α (0 + n)}'dir, \lc{w}'nin
        ise \lc{Vek α n}. Ya tipi \lc{Vek α (n + m)} olarak değiştirin
        (hangi durumda sorun çıkıyor?) ya da \lc{Nat.zero\_add n ▸ w}
        ile taşıyın.
  \item \lc{toList}'in taşımayla değişmediğini \lc{a15\_5}'teki gibi
        \leantac{subst} ile gösterin.
  \item Genelleştirilmiş önsav: \lc{(tersAux v w).toList = v.toList.reverse
        ++ w.toList}.
\end{enumerate}
\end{projekarti}

\begin{projekarti}{P11}{Tamsayıların halka yapısı}
\projebilgi{Bölüm~16}{\zorluk{2}}{2--3 gün}
\leandosya[P11]{Alistirmalar/Projeler}{p11}

\textbf{Önerilen yol.}
\begin{enumerate}[nosep]
  \item $(a - b)(c - d) = (ac + bd) - (ad + bc)$ olduğundan çarpmayı
        temsilciler üzerinde \lc{(a*c + b*d, a*d + b*c)} olarak tanımlayın.
  \item İyi tanımlılık doğrusal olmayan bir eşitlik ister; \leantac{omega}
        yetmez. \lc{Nat.mul\_add}, \lc{Nat.add\_mul} ile elle ya da Mathlib'in
        \leantac{ring}/\leantac{nlinarith} taktikleriyle ispatlayın.
        Önce bir argümanı, sonra diğerini değiştirmek işi kolaylaştırır.
  \item Halka aksiyomlarını \lc{Quotient.inductionOn₃} ile ispatlayın.
  \item \emph{Uzantı:} \lc{Tamsayi.isaret}'in bir halka eşyapı dönüşümü
        olduğunu ve \lc{Int}'e bir eşleme verdiğini gösterin.
\end{enumerate}
\end{projekarti}

\begin{projekarti}{P12}{Yapıcı model varlığı}
\projebilgi{Bölüm~10, 12}{\zorluk{3}}{1 hafta}
Bölüm~12'deki \lc{model\_varligi}, \lc{Classical.byContradiction} kullanıyordu.
Doğruluk tablosunu gerçekten tarayan yapıcı bir sürüm yazın.

\leandosya[P12]{Alistirmalar/Projeler}{p12}

\textbf{Önerilen yol.}
\begin{enumerate}[nosep]
  \item \lc{Γ}'nın formüllerindeki en büyük atom indisini bulun ve
        \lc{tumDegerlemeler} listesini \lc{List.find?} ile tarayın.
  \item \lc{p12\_dogru}, \lc{List.find?}'ın bulduğu elemanın koşulu
        sağladığını söyleyen önsavdan çıkar.
  \item \lc{p12\_tam} için P1'deki kapsama önsavı ve Bölüm~12'deki
        \lc{tamlik} gerekir.
  \item Son olarak \leankw{\#print} \lc{axioms} ile \lc{Classical.choice}'un
        kullanılmadığını doğrulayın.
\end{enumerate}
\end{projekarti}

% ---------------------------------------------------------------------
\section{Projeyi bitirdiğinizde}\label{sec:proje-bitis}

\begin{itemize}
  \item \textbf{\lc{sorry} kalmasın.} Projenizi \lc{Alistirmalar/}'dan
        kendi dosyanıza taşıyın ve sürekli tümleştirmenin \lc{sorry}
        denetimine dahil edin.
  \item \textbf{Aksiyomları listeleyin.} Sonucunuzun dayandığı aksiyomları
        belgeleyin. Yapıcı olması gereken bir sonuç \lc{Classical.choice}'a
        dayanıyorsa nedenini bulun.
  \item \textbf{Belgeleyin.} Her tanımın ve ana teoremin başına,
        neyi neden öyle tanımladığınızı açıklayan bir yorum yazın. Altı ay
        sonra kodu okuyacak kişi muhtemelen siz olacaksınız.
  \item \textbf{Paylaşın.} \Lean{} topluluğunun Zulip sohbet sunucusu ve
        Reservoir paket dizini, bitmiş projeler için doğal yerlerdir. Bir
        sonucun Mathlib'e katkı olarak uygun olup olmadığını orada
        sorabilirsiniz.
\end{itemize}

\begin{arastirmanotu}
Bu projelerin çoğu, ispat asistanlarının kıyaslama ölçütü olarak
kullanılmıştır. Church--Rosser ve güçlü normalleşme, POPLmark yarışmasının
\citep{aydemir2005} ruhuna uygun ``bağlayıcı içeren metateori''
örnekleridir. Kesme eleme ise formalize edilmiş ispat kuramının merkezindedir
\citep{chaudhuri2019,shillito2023}. Lean~4'te bu sonuçların bir kısmı
Foundation kütüphanesinde \citep{foundation} genel biçimleriyle bulunur.
Kendi çözümünüzü yazdıktan sonra oradaki çözümlerle karşılaştırmak,
tanım seçimlerinin etkisini görmenin en iyi yoludur.
\end{arastirmanotu}

% ---------------------------------------------------------------------
\chapter*{Sonsöz: Cilt II'ye Bakış}
\addcontentsline{toc}{chapter}{Sonsöz: Cilt II'ye Bakış}
\markboth{Sonsöz}{Sonsöz}

Cilt~I'e ``Mantık nedir?'' sorusuyla başladık ve \Lean{}'in çekirdeğinin
neye inandığını sorarak bitirdik. Bu yolda önermeler mantığını hem bir
araç hem de bir inceleme nesnesi olarak gördük. İspat sistemlerinin
geçerli ve tam olduğunu, kesmenin elenebildiğini ve tiplenebilen
programların ispat olduğunu \Lean{}'de ispatladık. Bölüm~16'da üçüncü hâlin
olanaksızlığının \Lean{}'de bir teorem olduğunu gördüğümüzde, Bölüm~4'te
açtığımız yapıcı/klasik sorusu da kapanmış oldu.

Cilt~I'in büyük ölçüde \emph{önermeler} mantığında kaldığını fark etmiş
olabilirsiniz. Niceleyicileri \Lean{}'in kendi mantığında bolca kullandık
(sığ gömme). Ama niceleyicili bir nesne dilini ve onun metateorisini derin
gömmeyle kurmadık. Bölüm~12'de Henkin yöntemini kâğıt üzerinde bırakmamızın
nedeni budur. Bağlı değişkenlerin bütün zorluğu, Bölüm~14'te
$\lambda$-hesabında gördüğümüz gibi, orada bizi bekliyor.

\textbf{Cilt~II} tam bu noktadan başlıyor. Planlanan içerik şöyle:

\begin{center}
\renewcommand{\arraystretch}{1.4}
\rowcolors{2}{tablosatir}{white}
\begin{tabularx}{0.92\linewidth}{>{\bfseries}p{0.7cm} X}
\rowcolor{tabloust}
\tablobaslik{} & \tablobaslik{Cilt II: İleri Konular} \\
1 & Birinci mertebe mantık: sözdizim, anlambilim ve Gödel--Henkin tamlığı, \Lean{}'de \\
2 & Kip mantığı ve Kripke anlambilimi; sezgici mantığın S4'e gömülmesi \\
3 & Kanıtlanabilirlik mantığı (GL) \\
4 & Gödel'in eksiklik teoremleri \\
5 & Küme kuramı, modeller ve bağımsızlık \\
6 & Karar verilebilirlik ve otomasyon: SAT/SMT, \leantac{omega}, \leantac{bv\_decide}, yansıma \\
7 & Metaprogramlama: kendi taktiğinizi yazmak \\
8 & Yapay zekâ ve otomatik teorem ispatı \\
9 & Büyük formalizasyon projeleri ve araştırma projeleri \\
\end{tabularx}
\end{center}

Cilt~I boyunca ``Cilt~II, Bölüm~$n$'' diye işaret ettiğimiz her konu bu
tablodadır. Cilt~I'deki P4 projesi (Kripke modelleri ile Peirce yasası) ve
P2 projesi (ispat aramanın doğruluğu), sırasıyla 2. ve 6. bölümlerin doğal
başlangıç noktalarıdır.

\vspace{1em}
\begin{flushright}
\yazar\\
\textit{Ekim 2026}
\end{flushright}
